% =====================================================================
%  CARNET D'ENTRAÎNEMENT — Fiche 12 (Chimie organique)
%  THÈME 3 — Nomenclature et priorité des fonctions
%  NIVEAU MOYEN — ÉNONCÉS
% =====================================================================
\documentclass[11pt,a4paper]{article}
\usepackage[utf8]{inputenc}
\usepackage[T1]{fontenc}
\usepackage{lmodern}
\usepackage[french]{babel}
\usepackage{amsmath,amssymb,mathtools}
\providecommand{\degree}{^\circ}
\usepackage{icomma}
\usepackage{siunitx}
\sisetup{output-decimal-marker={,},per-mode=symbol}
\usepackage[version=4]{mhchem}
\usepackage{chemfig}
\setchemfig{atom sep=2em,bond length=0.9em,cram width=2.5pt}
\usepackage{xcolor}
\usepackage{array}
\usepackage{booktabs}
\usepackage{tabularx}
\usepackage[most]{tcolorbox}
\usepackage{enumitem}
\usepackage{tikz}
\usetikzlibrary{arrows.meta,positioning,calc}
\usepackage{fancyhdr}
\usepackage{titlesec}
\usepackage[hidelinks]{hyperref}
\usepackage[a4paper,margin=2.1cm,headheight=26pt,headsep=9pt]{geometry}

\definecolor{primary}{RGB}{150,98,162}
\definecolor{primarydark}{RGB}{92,52,110}
\definecolor{excol}{RGB}{0,135,90}
\definecolor{attncol}{RGB}{200,40,55}
\definecolor{methcol}{RGB}{90,90,100}

\newcommand{\runtitle}{Thème 3 — Moyen — Énoncés}
\pagestyle{fancy}\fancyhf{}
\fancyhead[L]{\small\textcolor{primarydark}{\textbf{Fiche 12} — Chimie organique}}
\fancyhead[R]{\small\textcolor{primarydark}{\runtitle}}
\fancyfoot[C]{\small\thepage}
\renewcommand{\headrulewidth}{0.8pt}
\renewcommand{\headrule}{\hbox to\headwidth{\color{primary}\leaders\hrule height \headrulewidth\hfill}}
\setlength{\parindent}{0pt}\setlength{\parskip}{3pt}

\tcbset{boxbase/.style={enhanced,breakable,boxrule=0.9pt,arc=2.5pt,
   left=3mm,right=3mm,top=2mm,bottom=2mm,fonttitle=\bfseries,coltitle=white,
   toptitle=0.5mm,bottomtitle=0.5mm}}
\newtcolorbox[auto counter]{exo}[1]{boxbase,colback=primary!4,colframe=primary,
   title={\textsc{Exercice}~\thetcbcounter\ \textmd{--- \itshape #1}}}
\newtcolorbox{consigne}{boxbase,colback=methcol!8,colframe=methcol,sharp corners,
   title={\textsc{Consignes}}}

\newcommand{\banner}[2]{%
\begin{tcolorbox}[enhanced,colback=primary,colframe=primary,arc=5pt,boxrule=0pt,
   left=5mm,right=5mm,top=2.5mm,bottom=2.5mm,coltext=white]
{\large\bfseries #1}\\[1pt]{\small #2}
\end{tcolorbox}\vspace{2pt}}

\begin{document}

\banner{Thème 3 — Nomenclature et priorité des fonctions}{Niveau \textbf{Moyen} — Énoncés \quad$\vert$\quad Fiche 12, Chimie organique}

\begin{consigne}
Exercices \textbf{ouverts} : donne le \textbf{nom complet} en suivant la méthode (fonction principale $\to$ chaîne $\to$ numérotation $\to$ préfixes des fonctions secondaires par ordre alphabétique).
\end{consigne}

% ---------------------------------------------------------------------
\begin{exo}{acide + alcool : le préfixe hydroxy-}
Nomme la molécule ci-dessous. Précise la fonction principale, le suffixe et le préfixe utilisé pour la fonction secondaire.
\begin{center}\chemfig{HO-[:30]CH_2-[:-30]CH_2-[:30](=[:90]O)-[:-30]OH}\end{center}
\end{exo}

% ---------------------------------------------------------------------
\begin{exo}{cétone + halogène}
Nomme la molécule ci-dessous (repère la fonction principale, numérote pour lui donner le plus petit indice, puis place l'atome de chlore en préfixe) :
\begin{center}\chemfig{CH_3-[:30](=[:90]O)-[:-30]CH(-[:-90]Cl)-[:30]CH_3}\end{center}
\end{exo}

% ---------------------------------------------------------------------
\begin{exo}{aldéhyde + cétone : le préfixe oxo-}
La molécule ci-dessous porte \textbf{deux} carbonyles :
\begin{center}\chemfig{CH_3-[:30](=[:90]O)-[:-30]CH_2-[:30](=[:90]O)-[:-30]H}\end{center}
\begin{enumerate}[label=\textbf{\alph*)},leftmargin=2em,itemsep=1pt]
  \item Lequel des deux carbonyles est la fonction principale ? (utilise la priorité $\ce{-CHO} > \text{C=O}$)
  \item L'autre carbonyle (la cétone) devient un préfixe \textbf{oxo-}. Donne le nom complet.
\end{enumerate}
\end{exo}

% ---------------------------------------------------------------------
\begin{exo}{choisir le bon sens de numérotation}
On donne le pentan-2-ol :
\begin{center}\chemfig{CH_3-[:30]CH(-[:90]OH)-[:-30]CH_2-[:30]CH_2-[:-30]CH_3}\end{center}
\begin{enumerate}[label=\textbf{\alph*)},leftmargin=2em,itemsep=1pt]
  \item Numérote la chaîne dans les \textbf{deux} sens et indique, pour chacun, la position du groupe $\ce{-OH}$.
  \item Quel sens faut-il retenir, et pourquoi ? Donne le nom.
\end{enumerate}
\end{exo}

% ---------------------------------------------------------------------
\begin{exo}{aldéhyde + éther : le préfixe (alk)oxy-}
Nomme la molécule ci-dessous (l'éther, moins prioritaire que l'aldéhyde, se cite en préfixe « méthoxy- ») :
\begin{center}\chemfig{CH_3-[:30]O-[:-30]CH_2-[:30]CH_2-[:-30](=[:-90]O)-[:30]H}\end{center}
\end{exo}

\end{document}
